\section{Introduction}\label{sec:intro}

\subsection{Conifold transitions and their mirrors}
Let $X'$ be a projective threefold with trivial canonical bundle whose singular points $q\in Q$ are ordinary double
points, and let $\widehat X\to X'$ be a small resolution, with exceptional curves $C_q\cong\mathbb P^1$. A smoothing
$X_s$ of $X'$, when it exists, contains a vanishing three-sphere for each node, and passing from $\widehat X$ to $X_s$
through $X'$ is a conifold transition. The two sides are tied together by linear relations, a phenomenon that goes
back to Clemens \cite{Cle83}: in a projective conifold transition the relations among the vanishing spheres and those
among the exceptional curves are mutual annihilators \cite[Thm.~1.14]{LLW18}, so the more relations the curves
satisfy, the fewer the spheres satisfy. Under a hypothesis on the exceptional curves, Friedman identified the
image of the first-order deformations of $X'$ in the local deformations of the nodes with the relations
$\sum_q\lambda_q[C_q]=0$ in $H_2(\widehat X)$, up to a nonzero scalar at each node, and constructed smoothings from
relations with every coefficient nonzero \cite[Rem.~(4.5), Cor.~(4.7)]{Friedman}. With the unobstructedness theorems of Kawamata,
Ran and Tian this becomes a criterion, in the form proved by Namikawa: $X'$ is smoothable if and only if the exceptional
curves satisfy a relation with every coefficient $\lambda_q$ different from zero \cite[Thm.~2.5]{Nam02}; compare
\cite[Thm.~1.3]{RT09}.

Morrison proposed that mirror symmetry reverses conifold transitions \cite{Mor99}. If $\widehat X$ has a mirror family
$Y$, the exceptional curves of $\widehat X$ should correspond to vanishing spheres $\delta_q$ of a degeneration of $Y$,
and the degenerate members should have a small resolution mirror to the smoothing $X_s$. Morrison made the prediction
quantitative: if the relative Picard number of $\widehat X\to X'$ is $\rho_{\rm rel}$, the $|Q|$ nodes of the mirror should
impose only $\rho_{\rm rel}$ conditions on its complex moduli \cite[\S3.1]{Mor99}. In terms of lattices the prediction
reads as follows. Let $K\subset\Z^Q$ be the lattice of relations among the exceptional curves; its rank is
$|Q|-\rho_{\rm rel}$. On the mirror, the vanishing spheres should satisfy the same relations,
\begin{equation}\label{eq:intro-prediction}
\Bigl\{k\in\Z^Q:\ \sum_qk_q[\delta_q]=0\text{ in }H_3(Y;\Z)\Bigr\}=K,
\end{equation}
so that the members of the mirror family on which all the spheres $\delta_q$ collapse form a locus of codimension
$|Q|-\operatorname{rank}K=\rho_{\rm rel}$. Smith, Thomas and Yau proved the symplectic counterpart of Friedman's
criterion \cite[Thm.~2.9]{STY}: disjoint Lagrangian three-spheres admit a symplectic conifold transition, for one of the
$2^{|Q|}$ choices of small resolution, exactly when their classes satisfy a relation with every coefficient nonzero.
Batyrev and Kreuzer tested the picture on Batyrev hypersurfaces whose singular two-faces are unit squares
\cite{BK10}. They identified the curve classes with the circuits of the squares, specialised the mirror family to the
locus on which the face polynomial of every square factors into two binomials, where the local model at each square
suggests a conifold point, and observed that by \cite{STY} a projective small resolution of such a degenerate mirror
requires a relation among the vanishing spheres with every coefficient nonzero; they did not determine whether such a
relation exists \cite[\S3]{BK10}.

Casta\~no-Bernard and Matessi proposed a mechanism in terms of affine geometry \cite{CBM}. On an integral affine
three-manifold $B$ with singularities, a tropical $2$-cycle $S$ with boundary on the discriminant should produce, in a
small resolution of the conifold built from $B$, a $3$-chain bounded by a combination of exceptional curves, and in a
smoothing of the conifold built from the dual structure, a $4$-chain bounded by the same combination of vanishing
spheres. They carried this out for topological models glued from local pieces \cite[Thm.~7.3, \S9.5]{CBM} and
conjectured that the relations coming from tropical $2$-cycles exhaust the relations among the vanishing cycles
\cite[Conj.~8.3]{CBM}; for a family of compact examples they proved this conditionally on a basis conjecture
\cite[Prop.~9.2]{CBM}.

\subsection{The result}
This paper proves~\eqref{eq:intro-prediction}, with Morrison's count, for a large class of Batyrev hypersurfaces
and their mirror families, and constructs the relations from relative tropical $2$-cycles on the base of Gross's toric
degeneration.

The hypersurfaces are the following. Let $\Delta$ be a reflexive four-polytope and $X\subset\mathbb P_\Delta$ a
$\Delta^\circ$-regular anticanonical hypersurface. Singular two-faces of $\Delta$ give isolated singular points of $X$. We
assume that $\Delta$ is \emph{admissible} (Definition~\ref{def:admissible}): its edges are primitive, its singular
two-faces are unit squares, which give nodes, and lattice pentagons and hexagons with one interior lattice point and no
other lattice points besides their vertices, which give cones over del Pezzo surfaces of degrees seven and six, and the
edge of $\Delta^\circ$ dual to each of these faces has lattice length one, so that each face gives exactly one singular
point. Subdividing these faces into unit parallelograms and unimodular triangles, with a
choice of \emph{profile} $\sigma$ at each hexagon, defines a crepant partial resolution $X'_\sigma\to X$ whose only
singularities are nodes, one for each parallelogram \cite{Paper4,Paper5}. We assume that this partial resolution is
induced by a projective toric morphism; this is Hypothesis~\ref{hyp:P}. Then $X'_\sigma$ has a projective small
resolution $\widehat X$, and the relations among its exceptional curves in $H_2(\widehat X;\Z)$ are exactly the
lattice
\[
K=\Bigl\{\lambda\in\Z^Q:\ \sum_q\lambda_q\,\circv_q=0\Bigr\}
\]
of relations among the \emph{circuits} $\circv_q=e_b+e_d-e_a-e_c$ of the parallelograms $abcd$; the group
$H_2(\widehat X;\Z)$ is torsion-free (Proposition~\ref{prop:integral}). On the mirror side, $Y_t$
is a member of the Batyrev mirror family near its maximally degenerate limit $t\to0$, with real coefficients of a fixed
sign pattern, close to the locus where every parallelogram's face polynomial factors into two binomials. Near each
parallelogram the defining function of $Y_t$ has one nondegenerate critical point with small critical value, and
$\delta_q\subset Y_t$ is its vanishing sphere.

\begin{theorem*}[Theorem~\ref{thm:main}]
Under Hypothesis~\ref{hyp:P}, for $t$ small and coefficients in the window just described:
\begin{enumerate}
\item $\sum_qk_q[\delta_q]$ is a torsion class of $H_3(Y_t;\Z)$ only if $k\in K$, and it vanishes for $k\in K$. So the
relation lattice of the vanishing spheres is exactly $K$, and the $[\delta_q]$ span a free group $\Z^Q/K$ that meets the
torsion trivially.
\item Every relation $k\in K$ is the boundary of an explicit integral $4$-chain on $Y_t$ lifted from a relative
tropical $2$-cycle on the base of Gross's toric degeneration of $X'_\sigma$.
\item Near the given member, the members with a node at every one of these critical points form a nonempty complex
submanifold of codimension $|Q|-\operatorname{rank}K$, the relative Picard number of a projective small resolution
$\widehat X\to X'_\sigma$.
\end{enumerate}
\end{theorem*}

Item~(1) says that the vanishing spheres of the mirror satisfy exactly the same integral relations as the
exceptional curves of $\widehat X$, which is~\eqref{eq:intro-prediction}. Near the maximally degenerate limit the nodal locus of item~(3) lies close to, but in general not on, the locus
where the face polynomials factor; the terms of the defining function from the cells next to each parallelogram move
the critical values (Section~\ref{ssec:periodsremarks}).
Combined with the smoothing criterion, the theorem settles the
analogue of the question of Batyrev and Kreuzer for these families (Corollary~\ref{cor:smooth}): the vanishing spheres of the mirror satisfy a relation with every coefficient nonzero if
and only if $X'_\sigma$ is smoothable, and then one of the $2^{|Q|}$ conifold transitions of $Y_t$ in Lagrangian spheres
isotopic to the $\delta_q$ is symplectic.

The smallest example shows every feature. The polytope $\Delta_9$ of \cite{Paper5} has two unit squares and one
$dP_7$ pentagon, so $X'_\sigma$ has four nodes, and
$K=\Z\,(1,-1,-1,1)$. On the base $B$ of Gross's degeneration there is a disc whose boundary runs along the discriminant
through the four nodes, crossing each node straight, and whose node coefficients are $(1,-1,-1,1)$
(Section~\ref{ssec:x9}). Its lift is a $4$-chain on $Y_t$ bounded by $\delta_1-\delta_2-\delta_3+\delta_4$, and by the
theorem this is the only relation: the four spheres span a free group of rank three, and the members of the mirror family
with all four nodes form a locus of codimension three. Since $K$ contains a vector with no zero entry, $X'_\sigma$ is
smoothable, and the four spheres admit a symplectic conifold transition.

\subsection{Tropical 2-cycles}
The relations are produced on the base. Let $B=B'_\sigma$ be the dual intersection complex of Gross's toric degeneration
of $X'_\sigma$ \cite{Gross2005}: an integral affine three-sphere whose singularities lie along a graph $\Gamma$, the
discriminant. Each node $q$ gives a four-valent vertex $p_q$ of $\Gamma$, which is a negative node in the sense of
\cite{CBM} (Proposition~\ref{prop:base}); the other vertices of $\Gamma$ are trivalent of the two standard types. Each edge
of $\Gamma$ carries a vanishing vector, the direction fixed by the monodromy around it.

We work with \emph{relative tropical $2$-cycles} (Definition~\ref{def:relcycle}): integral $2$-chains on $B$ with
coefficients in the sheaf $\iota_*\Lambda$ of monodromy-invariant integral tangent vectors, whose boundary lies on
$\Gamma$ and runs along each edge in the direction of its vanishing vector. The boundary of such a chain is a
\emph{balanced network} on $\Gamma$, and at each node it has a \emph{node coefficient} $\rho_q\in\Z$, the jump of the
boundary coefficient across $q$ measured against the circuit $\circv_q$ (Definition~\ref{def:nodecoef}). The tropical
$2$-cycles of \cite{CBM} are examples; relative tropical $2$-cycles may in addition carry non-primitive fields, branch
arbitrarily, cross the discriminant with any invariant field, and pass a node in both directions. Their node
coefficients always lie in $K$, by an augmentation argument (Section~\ref{ssec:cycles-bdry}). The first main
theorem is that they fill it.

\begin{theorem*}[Theorem~\ref{thm:real}]
Every balanced network on $\Gamma$ bounds a relative tropical $2$-cycle, and the node coefficients of balanced networks
are exactly $K$. Hence the node coefficients of relative tropical $2$-cycles generate $K$ over $\Z$.
\end{theorem*}

The second main theorem lifts every relative tropical $2$-cycle $z$ to an integral $4$-chain $\check\Xi(z)$ on $Y_t$ with
\[
\partial\,\check\Xi(z)=\sum_q\rho_q(z)\,[\delta_q]
\]
(Theorem~\ref{thm:lift}). For a tropical $2$-cycle in the sense of \cite{CBM} the chain realises the four-dimensional
object that Casta\~no-Bernard and Matessi attach to it in a topological smoothing: away from the discriminant it is, up to
homotopy, the bundle of two-tori annihilated by the field $v$ of $S$ in the dual torus fibres.

On the resolution side we construct, for every tropical $2$-cycle $S$ in the sense of \cite{CBM}, an integral $3$-chain
on $\widehat X$ with boundary $\sum_q\rho_q(S)[C_q]$ (Theorem~\ref{thm:res}). So a single tropical object bounds on both
sides of the transition, with the same coefficients, as \cite{CBM} proposed. We conjecture that these strict tropical
$2$-cycles already generate $K$ (Conjecture~\ref{conj:strict}); this holds on $340$ of the $364$ hexagon-free polytopes
of our list with $K\neq0$ that satisfy~\ref{hyp:P}, and Theorem~\ref{thm:real} shows that the larger class of relative cycles
generates $K$ always.

\subsection{How the proof goes}
The proof of Theorem~\ref{thm:main} has three steps.

\emph{The base.} Theorem~\ref{thm:real} is combinatorial topology on $B\cong S^3$ with coefficients in $\iota_*\Lambda$.
At the boundary level, a balanced network is determined near each node by two integers, the flows through the two
pairs of opposite edges, and its node coefficient is their difference. Pushing the flows onto the two-faces of $\Delta$
turns the condition $\sum_q\lambda_q\circv_q=0$ into the vanishing of a boundary on the cell complex $\partial\Delta\cong
S^3$, and $H_1(S^3)=0$ produces a balanced network with node coefficients $\lambda$ for every $\lambda\in K$
(Section~\ref{ssec:cycles-bdry}). To show that every balanced network bounds, we push it off $\Gamma$ along the dual
cells and are left with a cycle in $B\setminus\Gamma$ with coefficients in the local system $\Lambda$ that must be shown
to bound modulo meridians of $\Gamma$. This is a complement lemma: a Mayer--Vietoris argument over a regular neighbourhood of $\Gamma$ shows
that $H_1(B\setminus\Gamma;\Z)$ is generated by meridians, and a computation on the boundary spheres of the facets of the
polytope whose boundary is $B$, where $H^1(S^2;\Z)=0$, shows that the coefficients cause no further obstruction. Every step is integral, so no torsion enters
(Section~\ref{sec:cycles}).

\emph{The lift.} The cycle lives on $B'_\sigma$, the base of $X'_\sigma$, while the mirror family degenerates over a second
base $B_Y$, built from the mirror polytope with a different subdivision. The cells of the two bases do not correspond,
so we first construct a piecewise linear homeomorphism $B'_\sigma\to B_Y$ that carries discriminant to discriminant and
the local system of integral tangent vectors to the dual local system, by a Legendre correspondence through a complex of
pairs of dual cells (Sections~\ref{sec:dualpair} and~\ref{sec:legendre}). On the mirror we replace $Y_t$ by a localized hypersurface $W^*$ in the
sense of Mikhalkin and Yamamoto \cite{Mikhalkin,Ya16}, in which each monomial is cut off where it is far from maximal; an
isotopy through smooth hypersurfaces carries $W^*$ to $Y_t$. The regime is chosen so that near each node $W^*$ is, in
holomorphic coordinates, the conifold smoothing $\{xy+z_1z_2=\kappa\}$ with $\kappa$ small and nonzero, whose vanishing
sphere is $\delta_q$ (Section~\ref{sec:regime}). The chain is then assembled from local pieces: thimbles at the nodes,
subsets of toric divisors along the boundary of the cycle, and two-torus bundles over the rest, extended over regions
of $W^*$ that are aspherical with fundamental group $\Lambda$. The lift is linear in the cycle, and we build it
constituent by constituent on an abstract domain, so that sheets that cross or coincide on the base simply add. The
only global obstruction is boundary in the class of the fibre three-torus, and it vanishes because the positive real
locus of $W^*$ meets each fibre torus once and misses every vanishing sphere (Section~\ref{sec:lift}).

\emph{Periods.} The first two parts give $K\subseteq\Rel(\delta)$. For the reverse inclusion we regard the periods
$\wp_q=\int_{\delta_q}\Omega$ of a holomorphic volume form as functions of the coefficients. Each extends holomorphically
over a neighbourhood of the given member, including the members with nodes, and equals the critical value
$\varepsilon_q$ times a function without zeros near the nodal locus, while $d\varepsilon_q$ is close to the circuit row of $q$. A relation
among the $[\delta_q]$ is therefore a linear relation among differentials whose matrix has rank
$|Q|-\operatorname{rank}K$, so $\Rel(\delta)$ has rank at most $\operatorname{rank}K$; since $K$ is saturated, equality
follows, and the same argument shows that no nonzero combination outside $K$ is torsion. The loci $\{\varepsilon_q=0\}$ for
$q$ in a basis of the row space meet transversally, and the remaining critical values vanish on their intersection;
this is item~(3) (Section~\ref{sec:periods}). The period argument is classical \cite{GMS95,LLW18}; what is needed
here is its uniform form near the maximally degenerate limit.

\subsection{Hypotheses}
Hypothesis~\ref{hyp:P} is the only restriction beyond admissibility. It is required by the method, since without it
the partial resolution $X'_\sigma$ is not projective over $X$ and has no Gross degeneration, and it is a genuine
restriction: on our list it fails for $11$ of the $375$ hexagon-free polytopes with $K\neq0$, and for every profile on
$59$ of the $142$ polytopes with a hexagon (Section~\ref{ssec:scope}). Two further conditions enter the construction and
hold automatically. The base requires a regular unimodular triangulation of $\partial\Delta^\circ$ with a height
satisfying Gross's convexity conditions; Theorem~\ref{thm:U} shows that one exists for every admissible polytope, by
a computer-assisted verification over the $51{,}827$ admissible polytopes of the Kreuzer--Skarke classification with
exactly checked certificates (Appendix~\ref{app:unimodular}). The mirror side requires a subdivision of $\partial\Delta$
in which the cells next to each node parallelogram have their other vertices at lattice height one; the pulling
refinement of the subdivision given by~\ref{hyp:P} has this property (Proposition~\ref{prop:pulling}).

\subsection{Relation to other work}
Casta\~no-Bernard and Matessi constructed Lagrangian torus fibrations over integral affine three-manifolds with
singularities \cite{CBM09} and the conifolds and cycles described above \cite{CBM}; the present paper realises their
construction on actual hypersurfaces and their Batyrev mirrors, where the base is Gross's \cite{Gross2005} and the
spaces are algebraic. Gross and Siebert developed the correspondence between toric degenerations and their affine data, with the discrete
Legendre transform that exchanges the two sides of the mirror \cite{GSI}; Haase and Zharkov
constructed the integral affine spheres of toric hypersurfaces and their torus fibrations \cite{HZ1}. Cycles in
Calabi--Yau hypersurfaces near a maximally degenerate limit have been built from tropical data in other settings:
Lagrangian rational homology spheres from tropical curves with ends on the discriminant by Mak and Ruddat
\cite{MR20}, and lifts of closed tropical $(p,q)$-cycles in hypersurfaces with unimodular tropicalisation by
Yamamoto \cite{Ya26}, who notes that the tropical $2$-cycles of \cite{CBM} can be regarded as tropical $(1,2)$-cycles
\cite[\S1, footnote]{Ya26}. Ruddat and Siebert compute period integrals over lifts of tropical one-cycles \cite{RS20},
Ruddat develops a homology theory of tropical cycles on integral affine manifolds with a perfect pairing \cite{Ru21},
and Abouzaid, Ganatra, Iritani and Sheridan compute periods tropically in the setting of the Gamma conjecture
\cite{AGIS20}. The chains constructed here differ from these in two respects: they are relative, with boundary at the
nodes, and the subdivision that defines the ambient toric variety of $X'_\sigma$ is not unimodular, since it contains the
node parallelograms. The period argument of Section~\ref{sec:periods} goes back to Greene, Morrison and Strominger
\cite{GMS95}; compare \cite{LLW18}.

\subsection{Organisation}
Section~\ref{sec:statements} fixes the setting and states the results. Part~I concerns the base:
Section~\ref{sec:base} describes the partial resolution, the relation lattice, Gross's base, its discriminant and the
mirror subdivision; Section~\ref{sec:dualpair} introduces the dual-pair complex, a cell structure shared by the base of
$X'_\sigma$ and the base of the mirror family; and Section~\ref{sec:cycles} proves Theorem~\ref{thm:real}. Part~II
concerns the mirror: Section~\ref{sec:legendre} transfers relative tropical $2$-cycles to the mirror base,
Section~\ref{sec:regime} constructs the localized mirror and the Morse charts at the nodes, Section~\ref{sec:lift} the
lift of relative tropical $2$-cycles (Theorem~\ref{thm:lift}), and Section~\ref{sec:periods} proves
Theorem~\ref{thm:main} and Corollary~\ref{cor:smooth}. Part~III, Section~\ref{sec:resolution}, constructs the chains on
the small resolution (Theorem~\ref{thm:res}); it is not used in Parts~I and~II. Part~IV, Section~\ref{sec:examples},
treats the examples and the list of polytopes, including the evidence for Conjecture~\ref{conj:strict}.
Appendix~\ref{app:unimodular} proves Theorem~\ref{thm:U}, and the remaining appendices contain deferred proofs and
data.

